% Vietnamese translation for OI Public Library
% Source-File: IOI中国国家候选队论文/2009/方展鹏/浅谈如何解决不平等博弈问题.ppt
\documentclass[11pt]{article}
\usepackage{oipl-vn}

\title{Bàn về cách giải bài toán trò chơi bất đối xứng\\{\large Ghi chú theo slide}}
\author{Fang Zhanpeng\\Trường Trung học số 1 Trung Sơn, Quảng Đông}
\date{2009}

\begin{document}
\maketitle

\origtitle{Solving partizan game problems}
\sourcepath{IOI-China-Candidate-Team-Papers/2009/Fang-Zhanpeng/partizan-games-slides.ppt}

\section{Chủ đề}

Trong trò chơi bất đối xứng, hai người chơi không có cùng tập nước đi. Vì vậy
cách dùng hàm SG quen thuộc cho trò chơi công bằng không áp dụng trực tiếp.
Slide giới thiệu một hướng nhìn bằng surreal number rồi quay lại các bài cụ
thể trong thi lập trình.

\section{Ý tưởng chính}

Một trạng thái được xem qua hai tập lựa chọn: các nước đi của bên trái và các
nước đi của bên phải. Giá trị của trạng thái không còn chỉ là một số SG, mà
phản ánh ưu thế tương đối giữa hai bên. Với bài thi, mục tiêu thực tế là tìm
được tiêu chuẩn thắng thua hoặc chiến lược chuyển trạng thái đủ hiệu quả.

\section{Ví dụ trong slide}

\begin{itemize}
  \item \textit{Procrastination}: dùng phân tích trạng thái để nhận ra vùng
  thắng và vùng thua khi hai bên có quyền đi khác nhau.
  \item \textit{The Easy Chase}: mô hình hóa vị trí của hai người chơi và suy
  luận theo khoảng cách, lượt đi và lựa chọn hợp lệ của từng bên.
\end{itemize}

\section{Khi đọc cùng bài viết}

Bài viết đầy đủ giải thích nền tảng lý thuyết kỹ hơn. Bản slide nên được đọc
như phần định hướng: xác định quyền đi riêng của mỗi bên trước, sau đó mới
chọn công cụ phân tích. Việc ép bài toán vào mô hình SG đối xứng quá sớm dễ
làm mất thông tin quan trọng.

\end{document}
